% Page 029

\seite{29}

3.\ Möchte die Regierung den jungen jetzigen Lehrer\ob
n. Kreis Hannover\ob
b. Würde es in der Fremde Vor\ohb
wärtsgehen gestalten.

4.\ der Schluß bestand in den Wünschen, als Gott\ob
möge unseren Kaiser noch viele Jahre\ob
dem deutschen Volke erhalten. Diesen Wunsch\ob
fand in einem Gruß ihrem Ausdruck.

\pend
\abschnitt{Einquartierung}
\pstart

Im Monat Juni erhielt der Ort zweimal\ob
Einquartierung von Theilen der Artillerie, welche\ob
auf dem Schießplatz Elsenborn ihre Übungen hinter\ob
von demselben zurückkehrten. Die Krieger\ob
zeigten sich nicht unbefriedigt von den guten\ob
Quartieren, welche ihnen in den Gasthöfen\ob
angewiesen wurden.

Im Monat Juli erhielt der Ort 2 Bataillone\ob
8 Wochen lang eine Einquartierung.\ob
Die Mannschaften waren mit der Verpflegung nicht\ob
unzufrieden.

\pend
\abschnitt{Lehrerwechsel}
\pstart

Mit dem 1.\ Sept. wurde der Lehrer Bauer in gleicher Eigen\ohb
schaft auf die nachgefasste Schule zu Großkönigsdorf bei\ob
Köln versetzt. Um die hiesige Schule wurden der Schul\ohb
amtsbewerber Wüst vernimmt. Das war indes wohl\ob
beim 22 Pfennigen sehr. Der Ort und die Kinder mit\ob
seinen Pfarrgenossen wirksam auf für ihn in jetzt schließen\ob
finden, daß in schon am folgenden Morgen nicht
\pend
